% ===== PRÉAMBULE CANONIQUE — à coller VERBATIM en tête de CHAQUE .tex =====
\documentclass[11pt,a4paper]{article}
\usepackage[utf8]{inputenc}\usepackage[T1]{fontenc}\usepackage{lmodern}
\usepackage[french]{babel}
\usepackage{amsmath,amssymb,mathtools}\usepackage{icomma}
\usepackage[version=4]{mhchem}          % \ce{...} : équations & formules
\usepackage{siunitx}\sisetup{locale=FR} % \qty{}{}, \si{}, \ang{}
\usepackage{chemfig}                    % \chemfig{...} : structures développées
\usepackage{xcolor}\usepackage{colortbl}
\usepackage{tikz}\usepackage{pgfplots}\pgfplotsset{compat=1.18}
\usepackage[most]{tcolorbox}\usepackage{enumitem}\usepackage{array,booktabs}
\usepackage[a4paper,margin=2.2cm]{geometry}
\usetikzlibrary{arrows.meta,calc,angles,quotes,positioning,patterns,backgrounds,shapes.geometric}
\definecolor{primary}{RGB}{222,49,99}\definecolor{primarydark}{RGB}{150,22,55}
\definecolor{defcol}{RGB}{0,114,178}\definecolor{methcol}{RGB}{52,92,148}
\definecolor{excol}{RGB}{0,135,90}\definecolor{attncol}{RGB}{200,40,55}
\tcbset{boxbase/.style={enhanced,breakable,boxrule=0.9pt,arc=2.5pt,
  left=3mm,right=3mm,top=2.3mm,bottom=2.3mm,fonttitle=\bfseries,coltitle=white}}
% --- boîtes TUTORIEL (noms EXACTS, ne pas renommer) ---
\newtcolorbox{cle}[1][]{boxbase,colback=defcol!6,colframe=defcol,
  title={\textsc{À retenir}\ifx\relax#1\relax\relax\else\textmd{ : \itshape #1}\fi}}
\newtcolorbox{methode}[1][]{boxbase,colback=methcol!7,colframe=methcol,sharp corners,
  title={\textsc{Méthode}\ifx\relax#1\relax\relax\else\textmd{ --- \itshape #1}\fi}}
\newtcolorbox{exemplebox}[1][]{boxbase,colback=excol!5,colframe=excol,
  title={\textsc{Exemple}\ifx\relax#1\relax\relax\else\textmd{ : \itshape #1}\fi}}
\newtcolorbox{piege}[1][]{boxbase,colback=attncol!6,colframe=attncol,
  title={\textsc{Piège}\ifx\relax#1\relax\relax\else\textmd{ : \itshape #1}\fi}}
\newtcolorbox{remarque}[1][]{boxbase,colback=black!4,colframe=black!45,
  title={\textsc{Remarque}\ifx\relax#1\relax\relax\else\textmd{ : \itshape #1}\fi}}
% --- boîtes EXERCICE / CORRIGÉ (noms EXACTS, auto-comptées) ---
\newtcolorbox[auto counter]{exo}[1][]{boxbase,colback=primary!4,colframe=primary,
  title={\textsc{Exercice}~\thetcbcounter\ifx\relax#1\relax\relax\else\textmd{ --- \itshape #1}\fi}}
\newtcolorbox[auto counter]{corrige}[1][]{boxbase,colback=excol!4,colframe=excol,
  title={\textsc{Corrigé}~\thetcbcounter\ifx\relax#1\relax\relax\else\textmd{ --- \itshape #1}\fi}}
\newcommand{\R}{\mathbb{R}}\newcommand{\N}{\mathbb{N}}\newcommand{\Z}{\mathbb{Z}}\newcommand{\ds}{\displaystyle}
% (un \usepackage{fancyhdr}+en-tête est possible ; il est ignoré côté web)
% =========================================================================
\begin{document}
\sisetup{per-mode=symbol}

\begin{exo}[acide ou base selon Brønsted ?]
Pour chaque espèce, dis si, en présence d'eau, elle joue le rôle d'un \textbf{acide} (elle cède un \ce{H+}) ou
d'une \textbf{base} (elle capte un \ce{H+}).
\begin{enumerate}[label=\alph*)]
  \item \ce{HCl}
  \item \ce{NH3}
  \item \ce{OH-}
  \item \ce{CH3COOH}
\end{enumerate}
\end{exo}

\begin{exo}[trouver la base conjuguée]
Donne la \textbf{base conjuguée} de chaque acide (on retire un \ce{H+}).
\begin{enumerate}[label=\alph*)]
  \item \ce{HF}
  \item \ce{NH4+}
  \item \ce{H2CO3}
  \item \ce{H3O+}
\end{enumerate}
\end{exo}

\begin{exo}[trouver l'acide conjugué]
Donne l'\textbf{acide conjugué} de chaque base (on ajoute un \ce{H+}).
\begin{enumerate}[label=\alph*)]
  \item \ce{NH3}
  \item \ce{CO3^{2-}}
  \item \ce{H2O}
  \item \ce{NO2-}
\end{enumerate}
\end{exo}

\begin{exo}[écrire la demi-équation d'un couple]
Écris la demi-équation acide-base \textbf{en milieu aqueux} (avec \ce{H2O} et \ce{H3O+}) pour chaque couple.
\begin{enumerate}[label=\alph*)]
  \item \ce{HNO2/NO2-}
  \item \ce{HCOOH/HCOO-}
\end{enumerate}
\end{exo}

\begin{exo}[électrolyte fort ou faible ?]
Pour chaque espèce, indique si c'est un électrolyte \textbf{fort} (dissociation totale, flèche \ce{->}) ou
\textbf{faible} (dissociation partielle, flèche \ce{<=>}), puis écris son équation de dissociation dans l'eau.
\begin{enumerate}[label=\alph*)]
  \item \ce{HCl}
  \item \ce{CH3COOH}
  \item \ce{NaOH}
  \item \ce{HF}
\end{enumerate}
\end{exo}

\end{document}
